% test-008.tex -- comandos \spnewread, \spread, \spusep, \spdsym y \spmsym
%
% Compilar:  lualatex-dev test-008.tex
% Validar:   show-pdf-tags --xml test-008.pdf | rnv-wrapp
%            verapdf --flavour ua2 --format text test-008.pdf
%
% Cada caso es un parrafo numerado, para poder referirse a el al escuchar
% el PDF. Los comentarios "doc:" indican lo que documenta el manual; en el
% resto, la lectura se revisa escuchando. Estos comandos no analizan los
% simbolos: solo cambian la forma en que se verbalizan.
% \spnewread trabaja en modo texto y registra cada simbolo una sola vez
% (es global); los demas, en modo matematico. Las entradas invalidas (que
% detienen la compilacion) no van aqui, entre ellas \spread con un
% simbolo sin registrar y sin las llaves read-arg ni silent.
\DocumentMetadata
  {
    lang=es-CL,
    pdfversion=2.0,
    pdfstandard=ua-2,
    tagging=on,
    tagging-setup={math/setup={mathml-SE}},
  }
\documentclass{article}
\usepackage[spanish]{babel}
\usepackage{unicode-math}
\usepackage{spintent}
\usepackage{hyperref}
\hypersetup
  {
    pdfdisplaydoctitle = true,
    pdftitle           = {test-008: utilidades de simbolos},
  }
\pagestyle{empty}
\begin{document}

\section*{A. Comandos spnewread y spread: símbolos registrados}

% Registro de simbolos (global, una sola vez cada uno). No imprime nada.
\spnewread{\cong}{es-congruente-con}%
\spnewread{\sim}{es-semejante-a}%
\spnewread{\perp}{es-perpendicular-a}%
\spnewread{\approx}{es-aproximadamente-igual-a}%

% doc: «es congruente con»
1. Símbolo registrado: $\spread{\cong}$.

% doc: «es semejante a»
2. Otro símbolo registrado: $\spread{\sim}$.

% doc: «es perpendicular a»
3. Perpendicular: $\spread{\perp}$.

% doc: «es aproximadamente igual a»
4. Aproximadamente igual: $\spread{\approx}$.

% doc: «triángulo blanco apuntando hacia arriba a b c es congruente con triángulo blanco apuntando hacia arriba d e f»
5. En una fórmula: $\triangle ABC \spread{\cong} \triangle DEF$.

% doc: «a b es perpendicular a c d»
6. Perpendicularidad: $AB \spread{\perp} CD$.

% doc: «x es aproximadamente igual a 3»
7. Aproximación: $x \spread{\approx} 3$.

% doc: «triángulo blanco apuntando hacia arriba a b c es semejante a triángulo blanco apuntando hacia arriba d e f es semejante a triángulo blanco apuntando hacia arriba g h i»
8. Dos símbolos en una misma fórmula: $\triangle ABC \spread{\sim} \triangle DEF \spread{\sim} \triangle GHI$.

\section*{B. Llave read-arg}

% doc: «es paralela a»
9. Ejemplo del manual: $\spread[read-arg={es-paralela-a}]{\parallel}$.

% doc: «es equivalente a»
10. Sobrescribe un símbolo registrado: $\spread[read-arg={es-equivalente-a}]{\cong}$.

% doc: «es menor o igual que»
11. Otro comando sin registrar: $\spread[read-arg={es-menor-o-igual-que}]{\leq}$.

% doc: «a b es paralela a c d»
12. En una fórmula: $AB \spread[read-arg={es-paralela-a}]{\parallel} CD$.

% doc: «ángulo a b c»
13. Con un símbolo de ángulo: $\spread[read-arg={ángulo}]{\angle}ABC$.

\section*{C. Llave silent}

% doc: «»
14. Con true y un símbolo registrado: $\spread[silent=true]{\cong}$.

% doc: «es congruente con»
15. Con false explícito: $\spread[silent=false]{\cong}$.

% doc: «»
16. Con true y un símbolo sin registrar: $\spread[silent=true]{\parallel}$.

% doc: «a b c d»
17. Con true en una fórmula: $AB \spread[silent=true]{\parallel} CD$.

% doc: «»
18. Alternativa documentada con read-arg: $\spread[read-arg={_-}]{\parallel}$.

% doc: «»
19. Con silent y read-arg a la vez: $\spread[silent=true,read-arg={es-paralela-a}]{\parallel}$.

\section*{D. Comando spusep}

% doc: «se abren paréntesis»
20. Paréntesis que abre: $\spusep{(}$.

% doc: «se cierran paréntesis»
21. Paréntesis que cierra: $\spusep{)}$.

% doc: «abrir corchetes»
22. Corchete que abre: $\spusep{[}$.

% doc: «cerrar corchetes»
23. Corchete que cierra: $\spusep{]}$.

% doc: «abrir llaves»
24. Llave que abre: $\spusep{\{}$.

% doc: «cerrar llaves»
25. Llave que cierra: $\spusep{\}}$.

% doc: «se abren paréntesis a más b se cierran paréntesis»
26. Par de paréntesis: $\spusep{(} a + b \spusep{)}$.

% doc: «abrir corchetes a más b cerrar corchetes»
27. Par de corchetes: $\spusep{[} a + b \spusep{]}$.

% doc: «el conjunto a coma b»
28. Par de llaves: $\spusep{\{} a, b \spusep{\}}$.

% doc: «se abren paréntesis a más b se cierran paréntesis»
29. Con size none, el valor por defecto: $\spusep[size=none]{(} a + b \spusep[size=none]{)}$.

% doc: «se abren paréntesis a partido por b se cierran paréntesis»
30. Con size big sobre una fracción: $\spusep[size=big]{(} \frac{a}{b} \spusep[size=big]{)}$.

% doc: «se abren paréntesis a partido por b se cierran paréntesis»
31. Con size Big sobre una fracción: $\spusep[size=Big]{(} \frac{a}{b} \spusep[size=Big]{)}$.

% doc: «se abren paréntesis a partido por b se cierran paréntesis»
32. Con size bigg sobre una fracción: $\spusep[size=bigg]{(} \frac{a}{b} \spusep[size=bigg]{)}$.

% doc: «se abren paréntesis a partido por b se cierran paréntesis»
33. Con size Bigg sobre una fracción: $\spusep[size=Bigg]{(} \frac{a}{b} \spusep[size=Bigg]{)}$.

% doc: «abrir corchetes a más b cerrar corchetes»
34. Corchetes grandes: $\spusep[size=Big]{[} a + b \spusep[size=Big]{]}$.

% doc: «el conjunto a coma b»
35. Llaves grandes: $\spusep[size=bigg]{\{} a, b \spusep[size=bigg]{\}}$.

% doc: «a más b»
36. Con silent true: $\spusep[silent=true]{(} a + b \spusep[silent=true]{)}$.

% doc: «se abren paréntesis a más b se cierran paréntesis»
37. Con silent false explícito: $\spusep[silent=false]{(} a + b \spusep[silent=false]{)}$.

% doc: «a partido por b»
38. Con size y silent: $\spusep[size=Big,silent=true]{(} \frac{a}{b} \spusep[size=Big,silent=true]{)}$.

\section*{E. Comando spdsym}

% doc: «a dividido b»
39. Por defecto: $a \spdsym b$.

% doc: «a dividido b»
40. Con set-sym old-style: $a \spdsym[set-sym=old-style] b$.

% doc: «a dividido b»
41. Con set-sym two-dots: $a \spdsym[set-sym=two-dots] b$.

% doc: «a dividido b»
42. Con set-sym slash: $a \spdsym[set-sym=slash] b$.

% doc: «a dividido b»
43. Con read-sym dividido: $a \spdsym[read-sym=dividido] b$.

% doc: «a dividido por b»
44. Con read-sym dividido-por: $a \spdsym[read-sym=dividido-por] b$.

% doc: «a dividido entre b»
45. Con read-sym dividido-entre: $a \spdsym[read-sym=dividido-entre] b$.

% doc: «a dividido por b»
46. Con old-style y dividido-por: $a \spdsym[set-sym=old-style,read-sym=dividido-por] b$.

% doc: «a dividido entre b»
47. Con slash y dividido-entre: $a \spdsym[set-sym=slash,read-sym=dividido-entre] b$.

% doc: «a se lee dividido b»
48. Con prefix: $a \spdsym[prefix={se-lee}] b$.

% doc: «a dividido entre b»
49. Con suffix: $a \spdsym[suffix={entre}] b$.

% doc: «a se lee dividido entre b»
50. Con prefix y suffix: $a \spdsym[prefix={se-lee},suffix={entre}] b$.

% doc: «12 dividido 4 es igual a 3»
51. Con números: $12 \spdsym 4 = 3$.

\section*{F. Comando spmsym}

% doc: «a por b»
52. Por defecto: $a \spmsym b$.

% doc: «a por b»
53. Con set-sym cross: $a \spmsym[set-sym=cross] b$.

% doc: «a por b»
54. Con set-sym dot: $a \spmsym[set-sym=dot] b$.

% doc: «a por b»
55. Con set-sym asterisk: $a \spmsym[set-sym=asterisk] b$.

% doc: «a por b»
56. Con read-sym por: $a \spmsym[read-sym=por] b$.

% doc: «a multiplicado por b»
57. Con read-sym multiplicado-por: $a \spmsym[read-sym=multiplicado-por] b$.

% doc: «a veces b»
58. Con read-sym veces: $a \spmsym[read-sym=veces] b$.

% doc: «a multiplicado por b»
59. Con cross y multiplicado-por: $a \spmsym[set-sym=cross,read-sym=multiplicado-por] b$.

% doc: «a veces b»
60. Con asterisk y veces: $a \spmsym[set-sym=asterisk,read-sym=veces] b$.

% doc: «2 x»
61. Con not-print true: $2 \spmsym[not-print=true] x$.

% doc: «2 por x»
62. Con not-print false explícito: $2 \spmsym[not-print=false] x$.

% doc: «2 x»
63. Con not-print y read-sym: $2 \spmsym[not-print=true,read-sym=multiplicado-por] x$.

% doc: «a se lee por b»
64. Con prefix: $a \spmsym[prefix={se-lee}] b$.

% doc: «a por entre b»
65. Con suffix: $a \spmsym[suffix={entre}] b$.

% doc: «3 por 4 es igual a 12»
66. Con números: $3 \spmsym 4 = 12$.

\section*{G. Combinaciones}

% doc: «12 dividido 4 por 2»
67. División y multiplicación: $12 \spdsym 4 \spmsym 2$.

% doc: «se abren paréntesis 6 por 2 se cierran paréntesis dividido 3»
68. Con delimitadores: $\spusep{(} 6 \spmsym 2 \spusep{)} \spdsym 3$.

% doc: «se abren paréntesis a por b se cierran paréntesis es aproximadamente igual a c»
69. Con un símbolo registrado: $\spusep{(} a \spmsym b \spusep{)} \spread{\approx} c$.

% doc: «se abren paréntesis a por b se cierran paréntesis dividido c»
70. En fórmula destacada:
\[ \spusep[size=Big]{(} a \spmsym b \spusep[size=Big]{)} \spdsym c \].

\section*{H. Comando spnewread: espacios, guiones bajos y comandos propios}

% Requiere spnewread-sanitize.patch. La lectura se normaliza: los espacios
% y los guiones bajos pasan a guiones, y se quitan los del inicio y del final.
\newcommand{\xcuad}{x^2}%
\spnewread{\neq}{es distinto de}%
\spnewread{\geq}{es_mayor_o_igual_que}%
\spnewread{\subset}{_es-subconjunto-de}%
\spnewread{\in}{pertenece a}%
\spnewread{\xcuad}{x al cuadrado}%

% doc: «a es distinto de b»
71. Lectura registrada con espacios: $a \spread{\neq} b$.

% doc: «a es mayor o igual que b»
72. Lectura registrada con guiones bajos: $a \spread{\geq} b$.

% doc: «a es subconjunto de b»
73. Lectura registrada con un guion bajo al inicio: $A \spread{\subset} B$.

% doc: «x pertenece a a»
74. Otra lectura con espacios: $x \spread{\in} A$.

% doc: «x cuadrado»
75. Un comando definido por el usuario: $\spread{\xcuad}$.

% doc: «y es igual a x cuadrado más 1»
76. El comando definido por el usuario en una fórmula: $y = \spread{\xcuad} + 1$.

% doc: «x cuadrado»
77. El mismo comando con silent true: $\spread[silent=true]{\xcuad}$.

% doc: «x cuadrado»
78. El mismo comando con read-arg: $\spread[read-arg={equis al cuadrado}]{\xcuad}$.

\end{document}
